\usepackage{graphicx}
\usepackage{booktabs}
\usepackage{hyperref}
\usepackage{geometry}
\usepackage{eurosym}
\usepackage[english, spanish]{babel}
\usepackage{csquotes}
\usepackage{fontspec}
\setmainfont{Libertinus Serif}
\usepackage{xcolor}
\definecolor{citascolor}{RGB}{229,7,7}
\definecolor{linkscolor}{RGB}{0,90,180}
\definecolor{urlscolor}{RGB}{100,46,157}
\usepackage{hyperref}
\hypersetup{
    colorlinks=true,
    linkcolor= linkscolor,
    citecolor= citascolor,
    urlcolor= urlscolor
}
\usepackage{listings}
% Para el entorno listings
\definecolor{darkgraycode}{RGB}{30,30,30}        % Fondo
\definecolor{lightgraytext}{RGB}{220,220,220}    % Texto base
\definecolor{keywordblue}{RGB}{86,156,214}       % Palabras clave
\definecolor{stringyellow}{RGB}{214,157,133}     % Cadenas
\definecolor{commentgreen}{RGB}{106,153,85}      % Comentarios
\definecolor{functionlightyellow}{RGB}{220,220,170} % Funciones
\definecolor{numbercolor}{RGB}{181,206,168}      % Números
\definecolor{decoratorpink}{RGB}{197,134,192}    % Decoradores
% Configuración de listings
\lstset{
  language=Python,
  basicstyle=\ttfamily\scriptsize\color{lightgraytext},
  keywordstyle=\color{keywordblue}\bfseries,
  commentstyle=\color{commentgreen}\itshape,
  stringstyle=\color{stringyellow},
  identifierstyle=\color{functionlightyellow},
  numberstyle=\tiny\color{lightgraytext},
  numbers=left,
  stepnumber=1,
  breaklines=true,
  showstringspaces=false,
  backgroundcolor=\color{darkgraycode},
  frame=single,
  rulecolor=\color{darkgraycode},
  captionpos=b,
  literate=
    {0}{{{\color{numbercolor}0}}}1
    {1}{{{\color{numbercolor}1}}}1
    {2}{{{\color{numbercolor}2}}}1
    {3}{{{\color{numbercolor}3}}}1
    {4}{{{\color{numbercolor}4}}}1
    {5}{{{\color{numbercolor}5}}}1
    {6}{{{\color{numbercolor}6}}}1
    {7}{{{\color{numbercolor}7}}}1
    {8}{{{\color{numbercolor}8}}}1
    {9}{{{\color{numbercolor}9}}}1
}
\usepackage{amsmath}
\usepackage[spanish]{cleveref}
\usepackage{amsfonts}
\usepackage{amssymb}
\usepackage{amsthm}
\usepackage[mathscr]{euscript}
\usepackage{mathrsfs}
\usepackage{textcomp}
\usepackage[
    sorting=none,
    style=ieee,
]{biblatex}
\bibliography{referencias}
\usepackage{float}
\usepackage{setspace}
\usepackage{tocloft}
\usepackage{caption}
\usepackage{subcaption}
\usepackage{longtable}
\usepackage{url}
\usepackage{fancyhdr}

\setlength{\parskip}{1em}
\geometry{
  top=1.75cm,
  bottom=2.5cm,
  left=2.0cm,
  right=2.0cm,
  headsep = 1cm
}

\setlength{\headheight}{42pt}
\pagestyle{fancy}
\fancyhf{} 

\fancyhead[L]{
  \includegraphics[height=1.3cm]{figuras/logo_UC.jpg}
}

\fancyhead[R]{
  \includegraphics[height=1.3cm]{figuras/logo_Evenbytes.png}
}

\fancyfoot[C]{\thepage}

% Mostramos y enumeramos hasta subsubsection
\setcounter{tocdepth}{3}
\setcounter{secnumdepth}{3}